\begin{pastexam}[label={px:ntut-115-ee-ds-24}]{北科大電機所（戊組）115 學年度 問答題第 4 題}
（問答）

Given the following C++ function, what will this function return for \textbf{abc\#\#d}. Explain your
reasoning. \textbf{(5\%)}
\begin{minted}[fontsize=\small]{cpp}
string processString(string input) {
    stack<char> S;
    string result = "";
    for(char c : input) {
        if(c == '#') { if(!S.empty()) S.pop(); }
        else { S.push(c); }
    }
    while(!S.empty()) { result = S.top() + result; S.pop(); }
    return result;
}
\end{minted}
\end{pastexam}

\begin{pastsoln}
答案：回傳 \texttt{"ad"}。

這個函式模擬「\texttt{\#} 是退格鍵」：一般字元 push 進堆疊，碰到 \texttt{\#} 就 pop 掉最近打的那個字元（堆疊是空的就
不做事，像在行首按退格）。 逐字追蹤（堆疊由底到頂寫）：
\begin{center}
\begin{tabular}{cll}
\toprule
字元 & 動作 & 堆疊 \\
\midrule
a & push & a \\
b & push & a b \\
c & push & a b c \\
\# & pop（刪掉 c） & a b \\
\# & pop（刪掉 b） & a \\
d & push & a d \\
\bottomrule
\end{tabular}
\end{center}
最後的 \texttt{while} 從頂端一個一個 pop，每次把頂端字元接在 \texttt{result} 的\textbf{前面}（\texttt{S.top() + result}）：
先 pop 出 \texttt{d}，\texttt{result} 是 \texttt{"d"}；再 pop 出 \texttt{a}，接在前面變成 \texttt{"ad"}。 因為 pop 的順序和放進去
的順序相反，接在前面剛好把順序轉回來。 如果寫成 \texttt{result + S.top()}，就會得到反過來的 \texttt{"da"}。

時間：每個字元最多 push 一次、pop 一次，迴圈是 $O(n)$；不過 \texttt{S.top() + result} 每次會產生新字串、複製 \texttt{result}，
最後那個迴圈實際上是 $O(n^2)$ 的字元複製，要省時間可以改成往後接再整個反轉。
\end{pastsoln}
